% GENERATED FILE. Edit canonical lessons, then run npm run generate:books.
% canonical-source-hash: fnv1a64:ff4fd0f4f9a683d4
% canonical-lessons: PA-C308-chara, PA-C308-goha, PA-C308-vachchha, PA-C308-sing, PA-C308-hakkna, PA-R308-review-302-308, PA-R308-review-final-308

\chapter{Feeding the Animals}
\label{ch:pa-a2-feeding-the-animals}
\hlchaptermodality{308}

\begin{chapteropening}
\textbf{By the end of this chapter:} \emph{I can talk about feeding and herding the animals.}

Complete the last lesson of chapter 308: I can talk about feeding and herding the animals.
\end{chapteropening}

\section[\texorpdfstring{\pa{ਚਾਰਾ} (chārā)}{ਚਾਰਾ (chārā)}]{\pa{ਚਾਰਾ} (chārā) — fodder}

\label{lesson:PA-C308-chara}



\begin{quote}
Before the new one: say the Punjabi for a teat, then the Punjabi for to graze.
\end{quote}

\subsection*{You'll want to know: \pa{ਚਾਰਾ}}
\textbf{\pa{ਚਾਰਾ}} — \emph{chārā} — \textquotedblleft{}fodder\textquotedblright{}.

\begin{quote}
\pa{ਮੈਂ} \pa{ਮੱਝ} \pa{ਲਈ} \pa{ਹਰਾ} \pa{ਚਾਰਾ} \pa{ਲਿਆਵਾਂਗਾ।}
\end{quote}

\begin{quote}
\emph{maiṁ majjh laī harā chārā liāvāṅgā.}
\end{quote}

\begin{quote}
I will bring green fodder for the buffalo.
\end{quote}

A masculine noun.

\begin{cousinweb}
\textbf{\pa{ਚਾਰਾ}} (\emph{chārā}) is inherited from Sanskrit \emph{cāra}, \textquotedblleft{}grazing, fodder\textquotedblright{}; CDIAL lists Punjabi \emph{cārā}, \textquotedblleft{}fodder\textquotedblright{}, under it (CDIAL 4755).
\end{cousinweb}

\begin{grammarlens}[title={What you've built}]
One more word for this chapter.
\end{grammarlens}

\subsection*{Your turn}
\begin{itemize}
\raggedright
  \item \emph{Say it:} \emph{chārā}
  \item \emph{Say it:} \emph{chārā}, once more
  \item \emph{Say it:} \emph{maiṁ majjh laī harā chārā liāvāṅgā.} — I will bring green fodder for the buffalo.
  \item \emph{Recall:} say the Punjabi for a teat, then the Punjabi for to graze, then say \emph{chārā} again
\end{itemize}

\begin{tcolorbox}[breakable,colback=teal!4,colframe=teal!35!black,title={Before you move on}]
What does \pa{ਚਾਰਾ} mean? (\textquotedblleft{}Fodder\textquotedblright{}.) Say it once more.
\end{tcolorbox}
\section[\texorpdfstring{\pa{ਗੋਹਾ} (gohā)}{ਗੋਹਾ (gohā)}]{\pa{ਗੋਹਾ} (gohā) — cow dung}

\label{lesson:PA-C308-goha}



\begin{quote}
Before the new one: say the Punjabi for to graze, then the Punjabi for fodder.
\end{quote}

\subsection*{You'll want to know: \pa{ਗੋਹਾ}}
\textbf{\pa{ਗੋਹਾ}} — \emph{gohā} — \textquotedblleft{}cow dung\textquotedblright{}.

\begin{quote}
\pa{ਗੋਹਾ} \pa{ਰੂੜੀ} \pa{ਉੱਤੇ} \pa{ਜਾਂਦਾ} \pa{ਹੈ।}
\end{quote}

\begin{quote}
\emph{gohā rūṛī utte jāndā hai.}
\end{quote}

\begin{quote}
The cow dung goes onto the manure heap.
\end{quote}

A masculine noun.

\begin{cousinweb}
\textbf{\pa{ਗੋਹਾ}} (\emph{gohā}) is inherited from Sanskrit \emph{gōśakṛt}, \textquotedblleft{}cowdung\textquotedblright{}; CDIAL lists Punjabi \emph{gohā}, \textquotedblleft{}cowdung\textquotedblright{}, under it (CDIAL 4333).
\end{cousinweb}

\begin{grammarlens}[title={What you've built}]
One more word for this chapter.
\end{grammarlens}

\subsection*{Your turn}
\begin{itemize}
\raggedright
  \item \emph{Say it:} \emph{gohā}
  \item \emph{Say it:} \emph{gohā}, once more
  \item \emph{Say it:} \emph{gohā rūṛī utte jāndā hai.} — The cow dung goes onto the manure heap.
  \item \emph{Recall:} say the Punjabi for to graze, then the Punjabi for fodder, then say \emph{gohā} again
  \item \emph{Write it:} copy \textbf{\pa{ਗੋਹਾ}} once, checking the whole word
\end{itemize}

\begin{tcolorbox}[breakable,colback=teal!4,colframe=teal!35!black,title={Before you move on}]
What does \pa{ਗੋਹਾ} mean? (\textquotedblleft{}Cow dung\textquotedblright{}.) Say it once more.
\end{tcolorbox}
\section[\texorpdfstring{\pa{ਵੱਛਾ} (vachchhā)}{ਵੱਛਾ (vachchhā)}]{\pa{ਵੱਛਾ} (vachchhā) — a calf}

\label{lesson:PA-C308-vachchha}



\begin{quote}
Before the new one: say the Punjabi for fodder, then the Punjabi for cow dung.
\end{quote}

\subsection*{You'll want to know: \pa{ਵੱਛਾ}}
\textbf{\pa{ਵੱਛਾ}} — \emph{vachchhā} — \textquotedblleft{}a calf\textquotedblright{}.

\begin{quote}
\pa{ਗਾਂ} \pa{ਦਾ} \pa{ਵੱਛਾ} \pa{ਖੇਤ} \pa{ਵਿੱਚ} \pa{ਦੌੜਦਾ} \pa{ਹੈ।}
\end{quote}

\begin{quote}
\emph{gāṁ dā vachchhā khet vich dauṛdā hai.}
\end{quote}

\begin{quote}
The cow's calf runs in the field.
\end{quote}

A masculine noun.

\begin{cousinweb}
\textbf{\pa{ਵੱਛਾ}} (\emph{vachchhā}) is inherited from Sanskrit \emph{vatsa}, \textquotedblleft{}calf, child\textquotedblright{}; CDIAL lists Punjabi \emph{vacchā}, \textquotedblleft{}suckling buffalo calf\textquotedblright{}, under it (CDIAL 11239).
\end{cousinweb}

\begin{grammarlens}[title={What you've built}]
One more word for this chapter.
\end{grammarlens}

\subsection*{Your turn}
\begin{itemize}
\raggedright
  \item \emph{Say it:} \emph{vachchhā}
  \item \emph{Say it:} \emph{vachchhā}, once more
  \item \emph{Say it:} \emph{gāṁ dā vachchhā khet vich dauṛdā hai.} — The cow's calf runs in the field.
  \item \emph{Recall:} say the Punjabi for fodder, then the Punjabi for cow dung, then say \emph{vachchhā} again
  \item \emph{Write it:} copy \textbf{\pa{ਵੱਛਾ}} once, checking the whole word
\end{itemize}

\begin{tcolorbox}[breakable,colback=teal!4,colframe=teal!35!black,title={Before you move on}]
What does \pa{ਵੱਛਾ} mean? (\textquotedblleft{}A calf\textquotedblright{}.) Say it once more.
\end{tcolorbox}
\section[\texorpdfstring{\pa{ਸਿੰਗ} (siṅg)}{ਸਿੰਗ (siṅg)}]{\pa{ਸਿੰਗ} (siṅg) — a horn}

\label{lesson:PA-C308-sing}



\begin{quote}
Before the new one: say the Punjabi for cow dung, then the Punjabi for a calf.
\end{quote}

\subsection*{You'll want to know: \pa{ਸਿੰਗ}}
\textbf{\pa{ਸਿੰਗ}} — \emph{siṅg} — \textquotedblleft{}a horn\textquotedblright{}.

\begin{quote}
\pa{ਬਲਦ} \pa{ਦੇ} \pa{ਸਿੰਗ} \pa{ਲੰਮੇ} \pa{ਹਨ।}
\end{quote}

\begin{quote}
\emph{balad de siṅg lamme han.}
\end{quote}

\begin{quote}
The ox's horns are long.
\end{quote}

A masculine noun.

\begin{cousinweb}
\textbf{\pa{ਸਿੰਗ}} (\emph{siṅg}) is inherited from Sanskrit \emph{śṛṅga}, \textquotedblleft{}horn\textquotedblright{}; CDIAL lists Punjabi \emph{siṅg} under it (CDIAL 12583).
\end{cousinweb}

\begin{grammarlens}[title={What you've built}]
One more word for this chapter.
\end{grammarlens}

\subsection*{Your turn}
\begin{itemize}
\raggedright
  \item \emph{Say it:} \emph{siṅg}
  \item \emph{Say it:} \emph{siṅg}, once more
  \item \emph{Say it:} \emph{balad de siṅg lamme han.} — The ox's horns are long.
  \item \emph{Recall:} say the Punjabi for cow dung, then the Punjabi for a calf, then say \emph{siṅg} again
  \item \emph{Write it:} copy \textbf{\pa{ਸਿੰਗ}} once, checking the whole word
\end{itemize}

\begin{tcolorbox}[breakable,colback=teal!4,colframe=teal!35!black,title={Before you move on}]
What does \pa{ਸਿੰਗ} mean? (\textquotedblleft{}A horn\textquotedblright{}.) Say it once more.
\end{tcolorbox}
\section[\texorpdfstring{\pa{ਹੱਕਣਾ} (hakkṇā)}{ਹੱਕਣਾ (hakkṇā)}]{\pa{ਹੱਕਣਾ} (hakkṇā) — to drive cattle; to herd}

\label{lesson:PA-C308-hakkna}



\begin{quote}
Before the new one: say the Punjabi for a calf, then the Punjabi for a horn.
\end{quote}

\subsection*{You'll want to know: \pa{ਹੱਕਣਾ}}
\textbf{\pa{ਹੱਕਣਾ}} — \emph{hakkṇā} — \textquotedblleft{}to drive cattle; to herd\textquotedblright{}.

\begin{quote}
\pa{ਸ਼ਾਮ} \pa{ਨੂੰ} \pa{ਮੈਂ} \pa{ਮੱਝਾਂ} \pa{ਘਰ} \pa{ਹੱਕਾਂਗਾ।}
\end{quote}

\begin{quote}
\emph{shām nū̃ maiṁ majjhāṁ ghar hakkāṅgā.}
\end{quote}

\begin{quote}
In the evening I will drive the buffaloes home.
\end{quote}

Past (I …): \pa{ਮੈਂ} \pa{ਮੱਝਾਂ} \pa{ਹੱਕੀਆਂ}, \emph{maiṁ majjhāṁ hakkīāṁ}.

\begin{cousinweb}
\textbf{\pa{ਹੱਕਣਾ}} (\emph{hakkṇā}) is inherited from Sanskrit \emph{hakkayati}, \textquotedblleft{}calls to\textquotedblright{}; CDIAL lists Punjabi \emph{hakkṇā}, \textquotedblleft{}to drive away\textquotedblright{}, under it (CDIAL 13939).
\end{cousinweb}

\begin{grammarlens}[title={What you've built}]
One more word for this chapter.
\end{grammarlens}

\subsection*{Your turn}
\begin{itemize}
\raggedright
  \item \emph{Say it:} \emph{hakkṇā}
  \item \emph{Say it:} \emph{hakkṇā}, once more
  \item \emph{Say it:} \emph{shām nū̃ maiṁ majjhāṁ ghar hakkāṅgā.} — In the evening I will drive the buffaloes home.
  \item \emph{Recall:} say the Punjabi for a calf, then the Punjabi for a horn, then say \emph{hakkṇā} again
  \item \emph{Write it:} copy \textbf{\pa{ਹੱਕਣਾ}} once, checking the whole word
\end{itemize}

\begin{tcolorbox}[breakable,colback=teal!4,colframe=teal!35!black,title={Before you move on}]
What does \pa{ਹੱਕਣਾ} mean? (\textquotedblleft{}To drive cattle; to herd\textquotedblright{}.) Say it once more.
\end{tcolorbox}
\section[(dialogue)]{Looking back: from Sowing Time to Feeding the Animals}

\label{lesson:PA-R308-review-302-308}



\begin{quote}
Say the Punjabi for a garden bed and a manure heap.
\end{quote}

\subsection*{Your turn}
Say these aloud:

\begin{itemize}
\raggedright
  \item a garden bed, a manure heap, to scatter, to sprout, mustard
  \item a canal, a well, to irrigate, a farmer, a Persian wheel
  \item a weeding hoe, to hoe, to ripen, an ear of grain, a grain
  \item reaping, to reap, a sickle, to thresh, a heap
  \item a winnowing basket, to winnow, chopped straw, rice straw, a sack
  \item a buffalo, a buffalo calf, to milk an animal, a teat, to graze
  \item fodder, cow dung, a calf, a horn, to drive cattle
\end{itemize}

\begin{tcolorbox}[breakable,colback=teal!4,colframe=teal!35!black,title={Before you move on}]
A garden bed? (\textbf{\pa{ਕਿਆਰੀ}}, \emph{kiārī}.) A manure heap? (\textbf{\pa{ਰੂੜੀ}}, \emph{rūṛī}.) To scatter? (\textbf{\pa{ਖਿਲਾਰਨਾ}}, \emph{khilārnā}.) To sprout? (\textbf{\pa{ਉੱਗਣਾ}}, \emph{uggṇā}.) Mustard? (\textbf{\pa{ਸਰ੍ਹੋਂ}}, \emph{sarhoṁ}.) A canal? (\textbf{\pa{ਨਹਿਰ}}, \emph{nahir}.) A well? (\textbf{\pa{ਖੂਹ}}, \emph{khūh}.) To irrigate? (\textbf{\pa{ਸਿੰਜਣਾ}}, \emph{sinjṇā}.) A farmer? (\textbf{\pa{ਕਿਸਾਨ}}, \emph{kisān}.) A Persian wheel? (\textbf{\pa{ਹਲਟ}}, \emph{halaṭ}.) A weeding hoe? (\textbf{\pa{ਖੁਰਪਾ}}, \emph{khurpā}.) To hoe? (\textbf{\pa{ਗੋਡਣਾ}}, \emph{goḍṇā}.) To ripen? (\textbf{\pa{ਪੱਕਣਾ}}, \emph{pakkṇā}.) An ear of grain? (\textbf{\pa{ਸਿੱਟਾ}}, \emph{siṭṭā}.) A grain? (\textbf{\pa{ਦਾਣਾ}}, \emph{dāṇā}.) Reaping? (\textbf{\pa{ਵਾਢੀ}}, \emph{vāḍhī}.) To reap? (\textbf{\pa{ਲੁਣਨਾ}}, \emph{luṇnā}.) A sickle? (\textbf{\pa{ਦਾਤੀ}}, \emph{dātī}.) To thresh? (\textbf{\pa{ਗਾਹੁਣਾ}}, \emph{gāhuṇā}.) A heap? (\textbf{\pa{ਢੇਰ}}, \emph{ḍher}.) A winnowing basket? (\textbf{\pa{ਛੱਜ}}, \emph{chhajj}.) To winnow? (\textbf{\pa{ਛੱਟਣਾ}}, \emph{chhaṭṭṇā}.) Chopped straw? (\textbf{\pa{ਤੂੜੀ}}, \emph{tūṛī}.) Rice straw? (\textbf{\pa{ਪਰਾਲੀ}}, \emph{parālī}.) A sack? (\textbf{\pa{ਬੋਰੀ}}, \emph{borī}.) A buffalo? (\textbf{\pa{ਮੱਝ}}, \emph{majjh}.) A buffalo calf? (\textbf{\pa{ਕੱਟਾ}}, \emph{kaṭṭā}.) To milk an animal? (\textbf{\pa{ਚੋਣਾ}}, \emph{choṇā}.) A teat? (\textbf{\pa{ਥਣ}}, \emph{thaṇ}.) To graze? (\textbf{\pa{ਚਰਨਾ}}, \emph{charnā}.) Fodder? (\textbf{\pa{ਚਾਰਾ}}, \emph{chārā}.) Cow dung? (\textbf{\pa{ਗੋਹਾ}}, \emph{gohā}.) A calf? (\textbf{\pa{ਵੱਛਾ}}, \emph{vachchhā}.) A horn? (\textbf{\pa{ਸਿੰਗ}}, \emph{siṅg}.) To drive cattle? (\textbf{\pa{ਹੱਕਣਾ}}, \emph{hakkṇā}.)
\end{tcolorbox}
\section[(dialogue)]{Feeding the Animals — once more}

\label{lesson:PA-R308-review-final-308}



\begin{quote}
Say the Punjabi for fodder and cow dung.
\end{quote}

\subsection*{Your turn}
\begin{itemize}
\raggedright
  \item \emph{Say it:} fodder, cow dung, a calf, a horn, to drive cattle
\end{itemize}

\begin{tcolorbox}[breakable,colback=teal!4,colframe=teal!35!black,title={Before you move on}]
Fodder? (\textbf{\pa{ਚਾਰਾ}}, \emph{chārā}.) Cow dung? (\textbf{\pa{ਗੋਹਾ}}, \emph{gohā}.) A calf? (\textbf{\pa{ਵੱਛਾ}}, \emph{vachchhā}.) A horn? (\textbf{\pa{ਸਿੰਗ}}, \emph{siṅg}.) To drive cattle? (\textbf{\pa{ਹੱਕਣਾ}}, \emph{hakkṇā}.)
\end{tcolorbox}
